\scanpage{7}
On est ramené à ceci

\uline{Corollaire} Si $n\geq2$, $\underline{\mathcal L}_n$ est isomorphe par les \unsure{hom induits} à $\underline{\mathcal L}_2^{\otimes n'}$ si $n$ pair, à $\underline{\mathcal L}_2^{\otimes n'}\otimes\underline{\mathcal L}_3$ si $n$ impair (où $n'\geq1$ dans le premier cas, $n'\geq0$ dans le deuxième).

Mais on a plutôt intérêt de procéder directement ainsi: On est ramené aussitôt par platitude coh. au cas où $S=\operatorname{Spec} k$, on peut alors un \unsure{élémt|élément} $f$ de $L_n$ non dans $L_{n-1}$, un $g$ de $L_m$ non dans $L_{m-1}$, donc $fg$ \deleted{\unsure{dans}} dans $L_{n+m}^{\unsure{\otimes}}$ mais manifestement pas dans $\unsure{L_{n-1}\otimes L_m|L_m\otimes L_n|L_{n+m-1}}$.

\uline{Corollaire} $\mathcal E_3$ engendre $\Gamma_*(f,\underline L)$

Considérons les opérations de $\underline{\mathfrak g}$ sur les germes de fonctions rationnelles sur $C$ [\unsure{restant|reste est} à définir], on constate que \deleted{$\underline{\mathfrak g}$} ces opérations induisent des hom
\[
\underline{\mathcal L}\otimes\underline{\mathcal L}_n\longrightarrow\underline{\mathcal L}_{n+1}\qquad n\geq2
\]

\uline{Proposition 3} \inserted{Soit $s\in S$, $p$ l'exp car. de $k(s)$.} L\deleted{es}'homomorphisme $\underline{\mathcal L}\otimes\underline{\mathcal L}_n\longrightarrow\underline{\mathcal L}_{n+1}$ est un \uline{isomorphisme} (\unsure{vu|en|voir} \deleted{\unsure{au|ordre}} sur $c$ \unsure{point|près}) ssi $p$ est premier à $n$

On est ramené au cas où $S=\operatorname{Spec} k$, et alors c'est une vérification locale: si $D$ est une \deleted{diff.} dérivation rég. non nulle en un pt\inserted{ $x\in C$}, et si $f\ne0$ sur $C$ a en $x$ un pôle d'ordre $n$,
% archive sheet number: 6
